% =====================================================================
%  courses/aritmetica/semana_01_logica/ejercicios.tex — ejercicios propuestos: conectivos lógicos y tabla de verdad
% ---------------------------------------------------------------------
%  Curso: Aritmética / Lógica | Semana: 01
%  MÓDULO: sin \documentclass ni \begin{document}
% =====================================================================
	
	\newpage
	\section{Ejercicios propuestos — Conectivos Lógicos y Tabla de Verdad}
	
	\begin{tcolorbox}[
		enhanced, colback=GrisClaro, colframe=AzulPizarra,
		arc=2pt, boxrule=0.7pt, left=8pt, right=8pt, top=4pt, bottom=4pt
		]
		\small\sffamily
		\textbf{Instrucciones:} Marca la alternativa correcta. Cada ejercicio tiene
		una única respuesta válida. Tiempo estimado: \textbf{35 minutos}.
		\hfill \textbf{Semana 01 — Conectivos Lógicos y Tabla de Verdad}
	\end{tcolorbox}
	
	\vspace{0.4cm}
	
	% --- NIVEL BÁSICO ----------------------------------------------------------
	\begin{tcolorbox}[
		enhanced, colback=VerdeSuave!40, colframe=VerdeOliva,
		arc=2pt, boxrule=0.5pt, fonttitle=\sffamily\bfseries\small,
		title={Nivel Básico}, left=4pt, right=4pt, top=3pt, bottom=3pt
		]
	\end{tcolorbox}
	
	\begin{multicols}{2}
		\setcounter{numejercicio}{0}
		
		% --- Ejercicio 1 -----------------------------------------------------------
		\ejercicio Señala la proposición \textbf{compuesta}:
		
		\begin{alternativas}
			\item Agripino y Cesarina son hermanos.
			\item Los Heraldos Negros es una obra de Vallejo.
			\item Joseph tomó la primera fotografía en blanco y negro.
			\item Carlos y Richard van juntos al cine.
			\item Daniel es profesor y Rosa es escritora.
		\end{alternativas}
		\vspace{0.3cm}
		
		% --- Ejercicio 2 -----------------------------------------------------------
		\ejercicio La conjunción $p \wedge q$ es \textbf{verdadera} cuando:
		
		\begin{alternativas}
			\item $p$ es V y $q$ es F
			\item $p$ es F y $q$ es V
			\item Ambas son F
			\item Ambas son V
			\item Al menos una es V
		\end{alternativas}
		\vspace{0.3cm}
		
		% --- Ejercicio 3 -----------------------------------------------------------
		\ejercicio El condicional $p \to q$ es \textbf{falso} únicamente cuando:
		
		\begin{alternativas}
			\item $p = F$ y $q = V$
			\item $p = F$ y $q = F$
			\item $p = V$ y $q = F$
			\item $p = V$ y $q = V$
			\item Ambas son falsas
		\end{alternativas}
		\vspace{0.3cm}
		
		% --- Ejercicio 4 -----------------------------------------------------------
		\ejercicio Una tabla de verdad con \textbf{3 variables} tiene:
		
		\begin{alternativas}
			\item 3 filas
			\item 4 filas
			\item 6 filas
			\item 8 filas
			\item 9 filas
		\end{alternativas}
		\vspace{0.3cm}
		
		% --- Ejercicio 5 -----------------------------------------------------------
		\ejercicio Si $p$ es V y $q$ es F, el valor de $p \vee q$ es:
		
		\begin{alternativas}
			\item F
			\item V
			\item Indeterminado
			\item Depende del contexto
			\item Ninguna de las anteriores
		\end{alternativas}
		
	\end{multicols}
	
	\vspace{0.3cm}
	
	% --- NIVEL INTERMEDIO ------------------------------------------------------
	\begin{tcolorbox}[
		enhanced, colback=AzulClaro!40, colframe=AzulMedio,
		arc=2pt, boxrule=0.5pt, fonttitle=\sffamily\bfseries\small,
		title={Nivel Intermedio}, left=4pt, right=4pt, top=3pt, bottom=3pt
		]
	\end{tcolorbox}
	
	\begin{multicols}{2}
		
		% --- Ejercicio 6 -----------------------------------------------------------
		\ejercicio La matriz principal de $(p \wedge q) \vee q$ es:
		
		\begin{alternativas}
			\item VVVV
			\item VFVF
			\item VVFV
			\item VVFF
			\item FFFF
		\end{alternativas}
		\vspace{0.3cm}
		
		% --- Ejercicio 7 -----------------------------------------------------------
		\ejercicio Si $\sim p \vee r$ es \textbf{falsa}, los valores de $p$ y $r$ son:
		
		\begin{alternativas}
			\item $p = V,\; r = V$
			\item $p = V,\; r = F$
			\item $p = F,\; r = V$
			\item $p = F,\; r = F$
			\item No puede ser falsa
		\end{alternativas}
		\vspace{0.3cm}
		
		% --- Ejercicio 8 -----------------------------------------------------------
		\ejercicio Simboliza: \textit{``Si tomas jugo de naranja o fresa, entonces
			estarás lleno''} ($p$: naranja, $q$: fresa, $r$: lleno):
		
		\begin{alternativas}
			\item $(p \wedge q) \to r$
			\item $(p \vee q) \to r$
			\item $p \to (q \vee r)$
			\item $(p \vee q) \leftrightarrow r$
			\item $r \to (p \vee q)$
		\end{alternativas}
		\vspace{0.3cm}
		
		% --- Ejercicio 9 -----------------------------------------------------------
		\ejercicio Para $p=V$, $q=F$, $r=V$, el valor de
		$(p \,\Delta\, q) \leftrightarrow \sim r$ es:
		
		\begin{alternativas}
			\item V
			\item F
			\item V y F alternan
			\item Indeterminado
			\item Depende de $s$
		\end{alternativas}
		\vspace{0.3cm}
		
		% --- Ejercicio 10 ----------------------------------------------------------
		\ejercicio Si $[(p \to q) \vee (q \vee \sim r)]$ es \textbf{falsa},
		el valor de $q$ es:
		
		\begin{alternativas}
			\item V
			\item F
			\item V o F indistintamente
			\item Solo si $p = V$
			\item No puede determinarse
		\end{alternativas}
		
	\end{multicols}
	
	\vspace{0.3cm}
	
	% --- NIVEL AVANZADO --------------------------------------------------------
	\begin{tcolorbox}[
		enhanced, colback=NaranjaSuave!60, colframe=NaranjaTierra,
		arc=2pt, boxrule=0.5pt, fonttitle=\sffamily\bfseries\small,
		title={Nivel Avanzado --- Tipo Admisión}, left=4pt, right=4pt, top=3pt, bottom=3pt
		]
	\end{tcolorbox}
	
	\begin{multicols}{2}
		
		% --- Ejercicio 11 ----------------------------------------------------------
		\ejercicio La proposición $(\sim p \wedge q) \to [(p \vee r) \vee t]$ es
		\textbf{falsa}. El valor de $\sim(\sim p \vee \sim q) \to (r \vee \sim t)$ es:
		
		\begin{alternativas}
			\item F
			\item V
			\item Contingente
			\item Tautología
			\item Contradicción
		\end{alternativas}
		\vspace{0.3cm}
		
		% --- Ejercicio 12 ----------------------------------------------------------
		\ejercicio Dada $\sim[(r \vee q) \to (r \to p)] \equiv V$ con $q$ falsa,
		el valor de $(r \vee \sim p) \wedge (q \vee p)$ es:
		
		\begin{alternativas}
			\item V
			\item F
			\item Depende de $r$
			\item Tautología
			\item Contradicción
		\end{alternativas}
		\vspace{0.3cm}
		
		% --- Ejercicio 13 ----------------------------------------------------------
		\ejercicio Si $(p \to \sim q) \vee (\sim r \to s)$ es \textbf{falsa},
		el valor de $(\sim p \wedge \sim q) \vee \sim p$ es:
		
		\begin{alternativas}
			\item F
			\item V
			\item No se puede determinar
			\item Igual al de $r$
			\item Igual al de $s$
		\end{alternativas}
		\vspace{0.3cm}
		
		% --- Ejercicio 14 ----------------------------------------------------------
		\ejercicio La proposición $(p \to q) \to \sim q$ es:
		
		\begin{alternativas}
			\item Tautología
			\item Contradicción
			\item Contingente
			\item Equivalente a $p$
			\item Equivalente a $\sim p$
		\end{alternativas}
		\vspace{0.3cm}
		
		% --- Ejercicio 15 ----------------------------------------------------------
		\ejercicio La proposición $[(\sim p \vee q) \to (q \leftrightarrow r)] \vee
		(q \wedge s)$ es \textbf{falsa} y $p$ es V. Los valores de $q$, $r$ y $s$
		(en ese orden) son:
		
		\begin{alternativas}
			\item V, V, F
			\item F, V, F
			\item F, F, F
			\item V, F, V
			\item F, V, V
		\end{alternativas}
		
	\end{multicols}
	
	% --- CLAVE DE RESPUESTAS ---------------------------------------------------
	\vspace{0.5cm}
	\begin{tcolorbox}[
		enhanced, colback=GrisClaro, colframe=GrisMedio,
		arc=2pt, boxrule=0.5pt,
		fonttitle=\sffamily\bfseries\small\color{GrisCarbón},
		title={Clave de Respuestas}
		]
		\small\sffamily
		\begin{multicols}{5}
			\noindent
			1.~E \quad 2.~D \quad 3.~C \quad 4.~D \quad 5.~B \\
			6.~B \quad 7.~B \quad 8.~B \quad 9.~B \quad 10.~B \\
			11.~B \quad 12.~B \quad 13.~B \quad 14.~C \quad 15.~C
		\end{multicols}
	\end{tcolorbox}
	
	% --- FIN — ejercicios.tex --------------------------------------------------